\documentclass{article}
\usepackage[utf8]{inputenc}
\usepackage[T1]{fontenc}
\usepackage{geometry}
\geometry{
    a4paper,
    left=2.5cm,
    right=2.5cm,
    top=2.5cm,
    bottom=2.5cm
}
\usepackage{graphicx}
\usepackage{float}
\usepackage{hyperref}
\usepackage[english]{babel}
\usepackage{amsmath}
\usepackage{amssymb}
\usepackage{array}
\usepackage{booktabs}
\usepackage{caption}
\usepackage{subcaption}
\usepackage{enumitem}
\usepackage{fancyhdr}
\usepackage{titlesec}
\usepackage{url}
\usepackage{textcomp}
\usepackage{xcolor}
\usepackage{listings}
\usepackage{multirow}
\usepackage{longtable}
\usepackage{tabularx}
\usepackage{times}

% Set up page style
\pagestyle{fancy}
\fancyhf{}
\rfoot{\thepage}
\renewcommand{\headrulewidth}{0pt}

% Title formatting
\titleformat{\chapter}[display]{\normalfont\bfseries}{\chaptertitle\thechapter}{0pt}{\Huge}
\titlespacing*{\chapter}{0pt}{50pt}{40pt}

% Document info
\title{Template Document}
\author{Academic Support Agent}
\date{\today}

\begin{document}
\maketitle

\frame{\titlepage}

\frame{



\section{Nombre de la secci n}

\begin{frame}{T tulo de la diapositiva}
\begin{itemize}
    \item Aqu  una ecuaci n:
\pause
\begin{equation*}
I_{a+}^{\alpha}f(x)=\frac{1}{\Gamma(\alpha)}\int_a^x(x-t)^{\alpha-1}f(t)dt.
\end{equation*}
\vspace{0.5 cm}
\pause
\item Aqu  otra ecuaci n con s mbolos previamente renombrados en el pre mbulo; tu puedes renombrar otros a tu gusto:
\begin{equation*}   \R,\Z,\Q,\C,\Sf^n,\N,\br\cdot,\cdot\ke
\end{equation*}
\end{itemize}
\end{frame}
\section{Nombre de la secci n}
\begin{frame}{T tulo de la diapositiva}
    \begin{itemize}
        \item Aqu  hago referencia a \cite{stein2010complex}.
        \pause\item Ahora un teorema:
        \pause\begin{teo}\label{teo_ejemplo}
            Si $A\in \R^{n\times n} ~(n\in\N)$ y $p_A$ es el polinomio caracter stico de $A$, entonces $p_A(A)=O$, siendo $O$ la matriz de ceros. 
        \end{teo}
    \end{itemize}
\end{frame}
\section{Nombre de la secci n}
\begin{frame}{Titulo de la diapositiva}
    \begin{figure}
        \centering
        % Content from Figura_ejemplo would be included here

        \caption{Ejemplo de una figura.}
        \label{fig_ejemplo}
    \end{figure}
\end{frame}
\section{Referencias}
\begin{frame}
\printbibliography    
\end{frame}
\end{document}